\documentclass[11pt,a4paper]{article}

\usepackage[margin=25mm]{geometry}
\usepackage[T1]{fontenc}
\usepackage{lmodern}
\usepackage[hidelinks]{hyperref}
% Shared, non-identifying metadata for the manuscript and title page.
\newcommand{\PaperTitle}{Completion-aware search with partial witness repair for dense fixed-shape facility layout}
\newcommand{\TargetJournal}{Computers \& Industrial Engineering}


\pagestyle{empty}
\setlength{\parindent}{0pt}
\setlength{\parskip}{0.65em}
\emergencystretch=1.5em

\begin{document}

\hfill \today

The Editors\\
\textit{\TargetJournal}

Dear Editors,

Please consider our manuscript, ``\PaperTitle,'' for publication as a research
article in \textit{\TargetJournal}.

Dense layout design calls for search decisions that preserve room for the
remaining facilities. Our paper uses a known feasible layout as reusable geometry
for subsequent construction. We ask when this information improves solution
quality under the same time budget, connecting constructive search with the use
of existing solutions in neighborhood design.

We develop a partial witness repair rule. After a proposed placement, the method
repacks only the displaced facilities and retains all unaffected rectangles in
the stored completion. Building on the feasible-continuation principle of
fortified rollout, this rule provides a targeted repair neighborhood for tightly
packed layouts. We keep the compatible completion separate from the best full
layout, and independently verify every candidate before it can update the
returned solution.

The main study compares partial repair, full completion regeneration, and a
validation-selected adaptive repair control on 240 attempted instances with
32 to 256 facilities, two geometry families, and fills of 90\% and 95\%.
The common initializer succeeds on 236 instances, yielding 2,124 runs.
At 95\% fill, partial repair exceeds the adaptive control by 3.45 to
5.91 percentage points in initial-objective improvement at 10 seconds and
1.68 to 4.12 points at 60 seconds. With 256 facilities, partial repair gives
higher quality than full regeneration at 10 seconds, while full regeneration
leads at 60 seconds. A separate one-second study isolates the benefit of
nonempty repair beyond unchanged witness reuse. Together, these results show
when targeted repair offers an early improvement and when broader
reconstruction benefits from additional time.

The paper addresses a computational decision problem in industrial engineering:
improving a feasible facility arrangement while preserving hard packing
constraints. Its contribution is a concrete repair mechanism and a controlled
evaluation of when retaining a feasible completion is worth its computational
cost. We believe this will interest readers of \textit{\TargetJournal} working
on facility layout, computer-aided planning and heuristic optimization.

The source code, benchmark generators, result files, and independent verification
scripts are publicly available at
\url{https://github.com/yeoneung/layout_optimization}. The Online Supplement
provides the complete budget comparisons and additional diagnostics.

% AUTHOR CONFIRMATION REQUIRED BEFORE SUBMISSION:
% Uncomment the paragraph below only after confirming both statements with
% all coauthors. Neither statement has been independently verified here.
% The manuscript is not under consideration by another journal. All authors
% have approved this manuscript and its submission to this journal.

Thank you for considering our manuscript.

Sincerely,\\
Yeoneung Kim\\
Corresponding author\\
Department of Industrial Engineering, Yonsei University\\
\href{mailto:yeoneung@yonsei.ac.kr}{yeoneung@yonsei.ac.kr}

\end{document}
